\item ¿Por qué es imposible la siguiente situación? Una barcaza transporta, a lo largo de un río, una carga de pequeños pedazos de hierro apilados. La pila de hierro tiene la forma de un cono cuyo radio $r$ de la base es igual a su altura central $h$. La barcaza tiene forma cuadrada, con lados verticales de longitud $2r$, así que la pila de hierro llena justo los bordes de la barcaza. El navío se aproxima a un puente de baja altura, y el capitán observa que la parte superior de la pila no logrará pasar por debajo del puente. El capitán ordena a la tripulación que tire al agua pedazos de hierro para así disminuir la altura de la pila. Conforme el hierro se va tirando, la pila siempre mantiene su forma cónica cuyo diámetro es igual a la longitud lateral de la barcaza. Después de un cierto volumen de hierro lanzado al agua, la nave pasa bajo el puente sin que la parte superior de la pila pegue en el puente.
